\section{Técnicas e instrumentos de recolección de datos}

A continuación, se detalla la correspondencia entre los métodos empíricos, las técnicas, los instrumentos y la finalidad de la recolección de datos durante el diagnóstico de campo:

\begin{table}[H]
\centering
\caption{Correspondencia entre métodos, técnicas, instrumentos y finalidad}
\label{tab:tecnicas}
\small
\begin{tabular}{p{2.5cm} p{3cm} p{3cm} p{5cm}}
\toprule
\textbf{Método Empírico} & \textbf{Técnica de Recolección} & \textbf{Instrumento Aplicado} & \textbf{Finalidad de la Recolección} \\
\midrule
Método Empírico (Cualitativo) & Entrevista semiestructurada & Guía de preguntas abiertas en profundidad & Comprender los procesos internos de atención, registro de datos y principales dificultades de comunicación en taller. \\
\addlinespace
Método Empírico (Cuantitativo) & Encuesta estructurada & Cuestionario estructurado (físico y Google Forms) & Cuantificar patrones de consulta, medios preferidos de comunicación y percepción de incertidumbre en clientes. \\
\addlinespace
Método Empírico (Documental) & Análisis y revisión documental & Ficha de registro y matriz de cotejo de comprobantes & Identificar los datos mínimos requeridos para la trazabilidad legal e histórica de las órdenes de servicio. \\
\bottomrule
\end{tabular}
\end{table}

\textbf{Procedimiento de Análisis de Datos:} Los datos cuantitativos recolectados mediante encuestas estructuradas se procesarán con estadística descriptiva (distribución de frecuencias y porcentajes) mediante hojas de cálculo (Excel), generando representaciones gráficas para caracterizar los patrones de consulta. Los datos cualitativos provenientes de las entrevistas semiestructuradas se analizarán mediante análisis temático de contenido, identificando categorías recurrentes sobre los problemas de registro y comunicación. La información resultante se integrará mediante triangulación metodológica para definir formalmente la especificación de requerimientos de software según la norma ISO/IEC/IEEE 29148:2018.

\textbf{Consideraciones Éticas:} La investigación se rige por principios de ética científica, garantizando la participación estrictamente voluntaria y el consentimiento informado previo de técnicos y clientes. Se aplica una política de minimización de datos y anonimización estricta, protegiendo la identidad de los participantes y el secreto comercial de los talleres en cumplimiento del Artículo 21 (inciso 2) y Artículo 130 de la Constitución Política del Estado de Bolivia y la Ley N.º 164.